%File: paper.tex
\documentclass[letterpaper]{article} % DO NOT CHANGE THIS
\usepackage[submission]{aaai2026}  % DO NOT CHANGE THIS
\usepackage{times}  % DO NOT CHANGE THIS
\usepackage{helvet}  % DO NOT CHANGE THIS
\usepackage{courier}  % DO NOT CHANGE THIS
\usepackage[hyphens]{url}  % DO NOT CHANGE THIS
\usepackage{graphicx} % DO NOT CHANGE THIS
\urlstyle{rm} % DO NOT CHANGE THIS
\def\UrlFont{\rm}  % DO NOT CHANGE THIS
\usepackage{natbib}  % DO NOT CHANGE THIS AND DO NOT ADD ANY OPTIONS TO IT
\usepackage{caption} % DO NOT CHANGE THIS AND DO NOT ADD ANY OPTIONS TO IT
\frenchspacing  % DO NOT CHANGE THIS
\setlength{\pdfpagewidth}{8.5in} % DO NOT CHANGE THIS
\setlength{\pdfpageheight}{11in} % DO NOT CHANGE THIS
\usepackage{xcolor}
\newcommand{\tk}[1]{\textcolor{blue}{[Tiankuo: #1]}}      % 你的批注，红色

% Permitted, commonly used packages.
\usepackage{amsmath}
\usepackage{amssymb}
\usepackage{booktabs}
\usepackage{subcaption}

\pdfinfo{
/TemplateVersion (2026.1)
}

\setcounter{secnumdepth}{1}

\title{Action Forcing}

\author{
    Anonymous Submission
    % Please add: Prof Zhong Fan (Z.Fan@exeter.ac.uk)
}
\affiliations{
    % Please acknowledge the Isambard project(s) where appropriate
    % Please acknowledge the EICDT where appropriate
}

\begin{document}

\maketitle

\begin{abstract}
We are living through a golden age of real-world video: good cameras, cheap storage, and an abundance of people filming content as far as the eye can see. However, corresponding egomotion labels remain scarce. Video diffusion models face no trouble learning from passive video but converting them into world models (WM) is dependent on finding good action labels. Data of this kind is rare, literature looks for other methods but these have proven hard to scale thus far. In this paper we present a method to recover an action space and corresponding actions from unsupervised video that scales with the data. 
We do this by making a simple action pipeline by stacking a pixel tracking method and using these pixel fields' principal components to derive the underlying egomotion.
While our method requires no action labels, it is data dependent and limited to recovering only the axes of egomotion present in the data. Because a high-capacity VDiTs can exploit the pixel-tracking objective, we train an online surrogate action critic using labels recovered by the pixel-space pipeline. The critic distils this signal into differentiable latent-space supervision while avoiding directly passing gradients through the decoder.
Finally we critique the use of video generation metrics to evaluate world models and introduce an example of an alternative, reference-free evaluation method. We measure \textit{controllability}, \textit{plausibility}, and \textit{geometric integrity}, revealing failures that conventional video metrics miss as rollouts can still remain sharp while relocating the scene or collapsing its structure.
Despite backwards actions being less than $1\%$ of our training data, we find that the model learns to reverse, scale its response linearly, and compose throttle with steering. Cross model action forcing shows that most evaluated systems either fail to reverse or lose the reverse component under diagonal commands. 
Ultimately, our results show that unlabelled real world video can be used to train WMs, removing the shackles of action labels from WM research and training.
\end{abstract}

\section{Introduction}
TODO:
0) Regenerate tables as score distributions
1) Combine all metrics to find legitimate rollouts
2) Stationary evaluations
3) Is zero motion a significant claim?
4) Take out half a page
5) Validate the no reference metrics
6) Put in generated videos
7) Check figure captions and figure references
8) Reference appendices in text
9) AI review

Large video diffusion transformers have recently found much success in learning powerful visual and temporal priors from internet-scale video. We naturally want to upgrade these models from video makers into world creators by conditioning their generation to respond to abstract egomotion inputs. World models have many potential applications such as training agents / humans in silico, intelligent planning, and expansive digital twins. However, the scale of training that unsupervised training bought LLM models are hard to achieve with world models as they currently need meticulously labelled data to function, data that simply does not exist naturally - new data either needs to be painstakingly collected \cite{navwormod, cameracontrol2, cosmos2, uniwm, xworld, gaia2}, or training is relegated to the domain of videogames: \cite{gamengen2024, diamond, olafworld, mineworld, gamefactory, oasis2} synthesised worlds that can only try to mimic the complexity and stochasticity of the real world. While some \cite{genie, gamengen2024, olafworld, adaworld} use latent action models to bridge the gap, these often prove to be complex models whose action spaces lack any meaningful grounding (cross context identifiability). Tellingly, even high performing systems retreat from messy data when they can: ~\cite{relic} abandons real world walking and drone data for unreal engine generated data, because real world data is dominated by forward motion and motion entangled across different dimensions. 

We train directly on this difficult regime. FrodoBots
video~\cite{frodobots} contains low-speed rover journeys through real cities,
with pedestrians, uneven ground, occlusion, weather, lens contamination, and
unstructured interactions. Its recorded commands are significantly affected by network latency and do not reliably describe realized motion; we find that the labels actively destabilise training. We therefore treat the videos as unlabelled and recover control from the motion rendered in the pixels.

Our approach rests on one primary idea - the pixel movements within frame follow the same broad stroke patterns for the same egomotion directions. Various actors present within the world move in their own ways and add complexity to the pixel movements but when this is taken over a large corpus of videos and the principal components of these pixel vectors are taken, the larger egomotion patterns emerge clearly. We use this as an action space as it ends up being linearly scalable, invertible, and composable with multiple directions. We use the two highest variance component weights as the egomotion input and we avoid conditioning on dense flow or raw tracks
directly~\cite{flovd,motionprompting,onlyflow}. This means we do not need to construct a full target motion
field at inference time.

We instantiate the method by adapting Wan2.1 on approximately $88$ hours of
FrodoBots video. Despite finding only approximately $55$ minutes of sustained
reverse motion in the full corpus, the resulting model learns backward
control and combines it well with steering.

Our contributions:
\begin{enumerate}

\item We recover an interpretable, signed, and composable egomotion action
    space directly from unlabelled real-world video using point tracking and
    PCA, without recorded actions or a learned latent-action model.

\item We introduce an online latent egomotion critic that distils
    pixel-space motion into differentiable latent supervision, grounding
    action following without backpropagating through the video decoder or
    tracker.

\item We show that standard video evaluation metrics are inadequate for
world models, where action forced rollouts have no paired reference to judge against. We introduce an example of a reference-free action-forcing evaluation that measures controllability, geometric integrity, and plausibility.

\end{enumerate}

\section{Related Work}

\paragraph{World models from pretrained video generators.}
Recent work adapts pretrained video DiTs into action-conditioned rollout
models using camera poses, keyboard--mouse inputs, game controls, or curated
robot actions~\cite{minwm,worldcam,astra,matrixgame2,hyworldplay,yume,
cosmos,dws,vid2world}. We compare directly against minWM, WorldCam, Astra,
Matrix-Game 2.0, HY-WorldPlay, and Yume. These systems demonstrate strong
controllability when actions are available, but rely on labelled,
instrumented, or simulated control signals. We instead recover a grounded
egomotion interface directly from unlabelled real-world video.

\paragraph{Actions recovered from video.}
Latent-action models learn controls through inverse dynamics without recorded
labels~\cite{genie,lapo,lapa,adaworld}, but their codes can be opaque and
inconsistent across contexts~\cite{olafworld}. Motion-conditioned generators
instead recover optical flow or dense point tracks from video
\cite{motionprompting,flovd,onlyflow}, but these representations are
high-dimensional and require a full motion field at inference time. We
compress tracked motion with PCA into a signed, low-dimensional action space
that supports direct commands, reversal, and composition.

\paragraph{VLM-based evaluation of generated video.}
Recent work increasingly uses multimodal language models as scalable proxies
for human evaluation of generated visual content. VIEScore uses
instruction-guided MLLMs to assess conditional image synthesis
\cite{viescore}, while VQAScore converts the probability of a
\emph{Yes} answer to a targeted visual question into an automatic metric
\cite{vqascore}. More directly, VideoPhy evaluates generated videos using
binary VLM judgments of semantic adherence and physical commonsense, and
constructs scores from normalized \emph{Yes}/\emph{No} token probabilities
\cite{videophy}. VideoScore learns a multi-aspect video evaluator from
human-rated generated videos \cite{videoscore}, and Video-Bench uses
structured MLLM queries across multiple dimensions of video-generation
quality \cite{videobench}. We follow this line of work, but use narrow,
reference-conditioned probes targeting style shift, introduced content, and
severe implausibility, and validate the resulting instruments against human
annotations.

\section{Preliminaries}
\subsection{Conditional Flow Matching}

Let $p_0 = \mathcal{N}(0,I)$ denote the base distribution, and let
$p_1(\cdot \mid c)$ denote the distribution of latent video chunks
conditioned on preceding frames $c$. Flow matching learns a
time-dependent vector field whose integration transports samples from
$p_0$ to the conditional data distribution:
\begin{equation}
\frac{d x_t}{d t} = v_\theta(x_t,t,c), \qquad t \in [0,1].
\end{equation}

Given a clean latent video chunk $x_1 \sim p_1(\cdot \mid c)$ and noise
$x_0 \sim \mathcal{N}(0,I)$, we use the linear probability path
\begin{equation}
x_t = (1-t)x_0 + t x_1,
\qquad
u_t(x_t \mid x_0,x_1) = x_1 - x_0.
\end{equation}

The model is trained by velocity regression:
\begin{equation}
\mathcal{L}_{\mathrm{flow}}
=
\mathbb{E}_{x_1 \sim p_1(\cdot \mid c)}
\left[
\left\|
v_\theta(x_t,t,c) - (x_1-x_0)
\right\|_2^2
\right].
\end{equation}

The target is determined by the sampled noise $x_0$ and clean latent
$x_1$. Diffusion time $t$ evolves the latent sample rather than
describing motion between video frames; temporal structure resides in
the latent video chunk and context $c$. We augment this velocity field
with the data-derived action introduced in \S~\ref{sec:planar}.
\section{Method}
\subsection{Action-Conditioned Flow}
\label{sec:planar}

A pretrained video DiT defines a context-conditioned velocity field
$v_\theta(x_t,t,c)$ whose integration produces a plausible continuation of
the preceding frames $c$. For brevity, let $z_t=(x_t,t,c)$ denote the current
latent state, flow time, and context. We expose a low-dimensional egomotion
interface to redirect the flow by augmenting this field with an action $a$:
\begin{equation}
    d x_t = v_\theta(z_t,a)\,dt.
    \label{eq:action_conditioned_ode}
\end{equation}
The context primarily determines the scene contents and plausible scene evolution while the action directs scene motion, effecting both generation and scene evolution. Because egomotion influences generation in this way, we condition the velocity network itself rather than applying a
fixed transformation to the completed sample.

At a common latent state, we isolate and ground the action induced change relative to a no-op command:
\begin{equation}
    \Delta v_\theta(z_t,a)
    =
    v_\theta(z_t,a)
    -
    v_\theta(z_t,0).
    \label{eq:action_flow_difference}
\end{equation}
The action representation determines the geometry that the conditional model
is asked to learn. Rover egomotion is continuous and low dimensional: zero
denotes no commanded motion, scalar changes represent command strength,
different control axes may act simultaneously, and changing the sign of an
axis requests the opposite motion. We therefore represent actions in a
centred affine coordinate system in which
\begin{equation}
    0,\qquad
    \lambda a,\qquad
    a_1+a_2,\qquad
    -a
\end{equation}
denote respectively no egomotion, scaled motion, compound motion, and motion
in the opposite direction. We use this structure as an inductive bias, asking the conditional flow to preserve the geometry of the command space without enforcing it as a hard architectural constraint. A generic
categorical or learned latent action representation can in principle learn
the same relationships, but the action space does not provide an intrinsic notion of neighbourhood, magnitude, addition, or inverse. Under finite data, those relationships must then be inferred separately. A centred affine chart
instead allows all commands to share the same scene dependent action
directions.

The visual realization of a command nevertheless remains scene dependent.
Depth, perspective, pose, visibility, occlusion, and scene geometry determine
how the same egomotion changes the video. We therefore do not assume a fixed
action-to-warp mapping. Instead, we hypothesize that, at a fixed $z_t$, the
action-induced change in the velocity field is locally control-affine:
\begin{equation}
    \Delta v_\theta(z_t,a)
    =
    G_\theta(z_t)a
    +
    \epsilon_\theta(z_t,a),
    \label{eq:implicit_control_affine}
\end{equation}
where $G_\theta$ is an implicit state-, time-, and scene-dependent action
operator, and $\epsilon_\theta$ captures action-dependent departures from
this affine approximation. The nonlinear realization of motion is contained
in $G_\theta$; the approximation concerns only how the low-dimensional
command selects and combines these scene-dependent responses.

When $\epsilon_\theta$ is small over the operational action range, the model
approximately transfers the geometry of the command space to its velocity
field:
\begin{equation}
    \begin{aligned}
        \Delta v_\theta(z_t,\lambda a_1+\gamma a_2)
        &\approx
        \lambda\Delta v_\theta(z_t,a_1)
        +
        \gamma\Delta v_\theta(z_t,a_2),\\
        \Delta v_\theta(z_t,-a)
        &\approx
        -\Delta v_\theta(z_t,a),\\
        v_\theta(z_t,0)
        &\approx
        v_\theta(z_t).
    \end{aligned}
    \label{eq:control_affine_predictions}
\end{equation}
The first relation captures smooth scaling and composition: command
magnitudes and combinations reuse common scene-dependent action directions
rather than defining separate conditional behaviours. The second captures
sign reversal: opposite egomotions occupy opposite sides of a shared axis.
This in theory permits abundant motion on one side of an axis (forward motion, perhaps) to inform its shared
response direction, while the scarcer opposite examples (backwards motion here) calibrate and test the negative branch. The
zero command provides the common reference from which every action-induced
change is measured. Actions are therefore represented as redirections of one
context-conditioned flow rather than as independent flow maps. Section~\ref{sec:pca}
constructs the centred action coordinates, while our response, composition,
and counterfactual evaluations test whether the learned generator transfers
their geometry to the flow.

\subsection{Data Derived Action Space}
\label{sec:pca}

Section~\ref{sec:planar} motivates a centred command space with consistent
notions of scale, composition, and sign. We recover such coordinates directly
from unlabeled video motion, while leaving the scene-dependent visual
realization of each command to the model.

In order to get fundamental movement signals, we use a pretrained off-the-shelf pixel tracking model \cite{cotracker} on a fixed $10\times10$ grid of starting pixels indexed by $j=1,\dots,100$. For every $0.5$s chunk, represented by 3 latent frames spanning 12 pixel frames, the tracker follows and returns the net displacement ($\Delta x, \Delta y,$) for each grid point across the chunk.
\begin{equation}
    m =
    \begin{bmatrix}
        \delta_1^\top &
        \delta_2^\top &
        \cdots &
        \delta_{100}^\top
    \end{bmatrix}^{\top}
    \in\mathbb{R}^{200}.
    \label{eq:movement_vector}
\end{equation}
The chunk is long enough for coherent displacement to exceed tracker noise,
but short enough that one approximately constant egomotion describes it and
successive commands may be composed autoregressively.

Let $\xi_{s,n}\in\mathbb{R}^2$ denote the unobserved planar egomotion of
chunk $n$ in scene $s$. Over a short interval, its stacked image-plane
displacement can be locally represented as
\begin{equation}
    m_{s,n}
    =
    J(c_{s,n})\,\xi_{s,n}
    +
    \rho_{s,n}
    +
    \varepsilon_{s,n},
    \label{eq:tracked_motion_model}
\end{equation}
where $J(c_{s,n})$ maps the rover egomotion into scene-dependent image
motion, $\rho_{s,n}$ contains motion independent of the rover, and
$\varepsilon_{s,n}$ contains tracking error and higher-order effects. The operator $J(c)$ accounts for depth, perspective, camera pose,
foreshortening, visibility, and occlusion, and may vary substantially between
scenes. Linearity is assumed only with respect to the short-horizon
egomotion $\xi$ at a fixed scene; image formation itself is not assumed to
be globally linear.

Rover egomotion produces coherent global motion patterns that recur
through the dataset. Independent pedestrians, vegetation, and tracker
errors are more localized and vary across scenes. Pooling observations
therefore emphasizes recurring egomotion structure while dispersing
scene-specific residual motion.

Let $\bar m=\mathbb{E}[m]$ denote the training-set mean. We compute
\begin{equation}
    \begin{aligned}
        C
        &=
        \mathbb{E}
        \left[
            (m-\bar m)(m-\bar m)^\top
        \right],\\
        C u_k
        &=
        \lambda_k u_k.
    \end{aligned}
    \label{eq:motion_pca}
\end{equation}

Although the scene-dependent operators $J(c)$ need not share an identical
column space, the rover has only two commanded planar degrees of freedom and
their global flow patterns recur across scenes. We therefore expect their
contribution to concentrate in a small leading eigenspace when it dominates
independent object motion and tracker noise. This is a data-dependent
assumption rather than a property guaranteed by PCA; we verify it through
the explained variance and coherent flow structure of the recovered
components.

Let $U=[u_1,\ldots,u_d]$ contain the retained PCA directions. We project each
tracked-motion vector into this basis as
\begin{equation}
    p = U^\top(m-\bar m).
    \label{eq:pca_projection}
\end{equation}
The resulting scores are scaled using their training-set variation and passed
through a $\tanh$ function to obtain bounded actions $a\in(-1,1)^d$. This
preserves the centred, signed ordering of each PCA direction while reducing
the influence of extreme motions. The resulting action space is locally
linear near its centre and smoothly saturates toward its boundaries.

The nonlinear squash may produce a small nonzero empirical mean in the
bounded action codes, but the coordinate origin remains
$a=\mathbf{0}$ and is used as the reference command at inference. The raw
PCA scores define an affine coordinate system in the retained tracked-motion
subspace: scaling changes represented motion proportionally, scores add to
form compound motion, and $-p$ denotes the opposite displacement. Because
the columns of $U$ are orthogonal, these coordinates are uniquely recoverable
within the retained subspace. The bounded actions preserve the origin,
ordering, and sign reversal, while proportional scaling and composition are
local approximations due to saturation. Whether the DiT transfers this
geometry to its conditional velocity field remains the empirical hypothesis
of Eq.~\eqref{eq:implicit_control_affine}.

On FrodoBots, the leading component captures forward/backward motion and
explains approximately $60\%$ (almost all forward) of tracked motion variance, while the second captures yaw and explains approximately $20\%$. Together, they align with the rover's two commanded planar degrees of freedom without using the recorded control labels. PCA leaves component signs arbitrary, so we fix each sign
once for interpretability and use that convention throughout training and
evaluation.

Genuine backward motion constitutes less than one hour of the
$2{,}000$ hour raw dataset, but occupies the negative side of the shared
throttle coordinate rather than requiring a separate action category.
Under the control-affine hypothesis, abundant forward examples constrain the
shared throttle response, while the sparse backward examples calibrate and
test its negative branch. Whether the generator successfully
inherits this sign structure is evaluated using counterfactual commands in
\S\ref{sec:eval}.

\subsection{Latent Egomotion Critic without Decoder Gradients}
\label{sec:critic}

The two-dimensional action space of \S\ref{sec:pca} provides the command
geometry required by \S\ref{sec:planar}, but flow matching alone does not
ensure that the generated video realizes the requested motion. We therefore
measure the egomotion rendered by the current generator and use it as an
additional training signal. The DiT is conditioned only on the leading two
PCA dimensions, corresponding to throttle and yaw, while an online critic
predicts the leading $r=8$ motion components to provide richer supervision.

Let $y\in\mathbb{R}^{8}$ denote the retained PCA motion code and let
$a=y_{1:2}\in\mathbb{R}^{2}$ contain its leading throttle and yaw
components. Only $a$ is supplied to the DiT; all eight components of $y$ are
used for training-time critic supervision.

Direct supervision would require decoding the generated latent, tracking its
pixels, projecting the tracks into the PCA basis, and differentiating through
the entire pipeline. In our setting this expands the loss interface from
$21\times16\times60\times104$ latent values to
$81\times3\times480\times832$ RGB values, approximately $46\times$ more
activations without even accounting for the decoder and tracker. More importantly, pixel-level gradients may reward local image displacement without requiring
coherent scene-level egomotion. We therefore use the decoder and tracker only
as a frozen, no-gradient visual teacher.

At each training step, the current velocity prediction is converted into a
one step clean latent estimate $\hat{x}_{\mathrm{clean}}$ in pixel space. We apply the visual teacher to this estimate at every sampled flow timestep, without
completing the generation trajectory or any time based filtering.

Let $\mathcal{T}_8$ denote the frozen pixel-space readout consisting of
latent decoding, point tracking, displacement assembly, PCA projection,
scaling, and squashing from \S\ref{sec:pca}. The motion rendered by the
current sample is measured as
\begin{equation}
    y^{\mathrm{teach}}
    =
    \mathcal{T}_8
    \left(
        \mathrm{sg}(\hat{x}_{\mathrm{clean}})
    \right)
    \in\mathbb{R}^{8},
    \label{eq:teacher_readout}
\end{equation}
where $\mathrm{sg}(\cdot)$ denotes stop-gradient. The supervisory label is
therefore obtained through a full pixel-space motion measurement, but this
pathway is forward-only: neither the decoder nor the tracker transmits
gradients to the generator.

A lightweight critic $D_\phi$ learns to predict the same motion code directly
from the same generated latent (in latent space):
\begin{equation}
    \hat y_\phi
    =
    D_\phi
    \left(
        \mathrm{sg}(\hat{x}_{\mathrm{clean}}),\tau,a
    \right).
    \label{eq:critic_prediction}
\end{equation}
The critic also receives the sampled flow timestep because the quality and
statistics of $\hat{x}_{\mathrm{clean}}$ vary with respect to diffusion time.
\begin{equation}
    \mathcal{L}_{\mathrm{crit}}
    =
    \left\|
        W_c
        \left(
            \hat y_\phi-y^{\mathrm{teach}}
        \right)
    \right\|_2^2.
    \label{eq:critic_loss}
\end{equation}
The first two dimensions receive full weight because they correspond to the
commanded throttle and yaw axes. The remaining six receive reduced weight,
providing a richer description of rendered motion without expanding the
two-dimensional control interface. Because the critic
is trained online on samples from the generator before it is asked to supervise those same latents, its approximation is fresh and continuously on manifold.
\begin{equation}
    \mathcal{L}_{\mathrm{act}}
    =
    \left\|
        W_g
        \left(
            D_{\bar\phi}(\hat{x}_{\mathrm{clean}},\tau)
            -
            y
        \right)
    \right\|_2^2.
    \label{eq:latent_action_loss}
\end{equation}
Here, $y$ is the eight-dimensional PCA motion code of the real training
chunk from which the conditioning action $a=y_{1:2}$ was obtained. The first
two dimensions enforce the requested throttle and yaw, while the lower-weight
remaining dimensions constrain the broader motion pattern rendered by the
sample.

The expensive motion label is thus computed in pixel space, while the
differentiable critic operates entirely in latent space. Because the critic
is trained online on samples from the evolving generator, its approximation
continually adapts to the latent distribution on which its gradients are
used.

The complete generator objective is
\begin{equation}
    \mathcal{L}_{\mathrm{gen}}
    =
    \mathcal{L}_{\mathrm{flow}}
    +
    \lambda_{\mathrm{act}}
    \mathcal{L}_{\mathrm{act}},
    \label{eq:generator_objective}
\end{equation}
while $\mathcal{L}_{\mathrm{crit}}$ updates only $\phi$. Thus,
\S\ref{sec:pca} supplies the visual action coordinates, the frozen teacher
measures whether generated samples realize them, and the latent critic
transfers that measurement into the action-conditioned flow of
\S\ref{sec:planar} without decoder or tracker gradients. The teacher and
critic are used only during training and introduce no inference-time cost.

We train and generate autoregressively at the level of temporal latent chunks.
Each training video contains seven chunks and we train using the teacher forcing paradigm \cite{williams1989learning}. We decompose the seven chunk sequence into seven teacher-forced prediction sequences of varying lengths, each targeting one chunk with a different
amount of preceding ground-truth context. Every target therefore has a clean
reference, and a single video provides seven supervised transitions.

At inference, we mirror this factorization by generating one latent chunk at
a time and appending it to the context before predicting the next. This avoids
denoising the complete rollout jointly, keeps the prediction target local, and
allows actions to be changed between successive chunks. When inference is run with 48 diffusion steps like in training, it takes 86s (on a GH200) to generate 7 chunks (or \~6 seconds) of content autoregressively, which at 16 FPS gives a generation speed of 1.11 FPS. Since Wan is a flow model, 20 diffusion steps suffice to generate with minimal loss in quality which subsequently increases the generation FPS to 2.72 FPS.  

\section{Evaluation}
\label{sec:eval}

\paragraph{Setup.}
We train on ${\sim}88$ hours of FrodoBots video comprising $2{,}577$ rides and ${\sim}61$K windows across eleven cities. Genuine reverse motion is exceptionally rare: the full ${\sim}2{,}000$-hour corpus contains only $632$ sustained reverse windows (${\sim}55$ minutes), which we oversample $6\times$. Training dynamics are analysed on four fixed training rides using both recorded commands (GT) and sign-flipped counterfactual commands (FLIP).

For cross-model evaluation we define an action-forcing protocol: 32 stationary contexts from held-out rides, each continued under eight constant compass commands {F, FR, R, BR, B, BL, L, FL}, giving 256 rollouts per model. Every model is rolled out over a common six-second horizon at its native frame rate and through its native autoregressive interface.

\paragraph{Comparisons.}
We compare against minWM~\cite{minwm}, WorldCam~\cite{worldcam},
Astra~\cite{astra}, Matrix-Game 2.0~\cite{matrixgame2},
HY-World 1.5 WorldPlay~\cite{hyworldplay}, and
Yume-1.5~\cite{yume}. Each compass command is translated into the model's native pose, keyboard–mouse, or text interface and held constant throughout the rollout. Because these interfaces differ in units and permitted command
strengths, we compare realized direction, forward--backward asymmetry, and
translation--steering composition rather than absolute control gain.
Complete command mappings are provided in
Appendix~\ref{app:baseline_controls}. We additionally ablate training batch
size, the number of PCA dimensions supervised by the critic, and the action
conditioning pathway.

\subsection{Controllability}

We measure controllability as agreement between the commanded action and the
motion rendered by the generated video, using the same frozen
CoTracker--PCA readout for every model. Because this readout also supplies our training signal, agreement with
it is not an independent success criterion for our model. However, we are not
aware of a more generally applicable alternative. Image-based estimators of
zoom, translation, or roll assume a rigid scene undergoing a projective
transformation, and become unreliable when generators warp three-dimensional
structure or when several motion axes must be evaluated jointly.

Our readout is deliberately simple. It tracks a fixed grid of pixels and
projects their dominant displacements onto a frozen PCA basis, with no learned
parameters or architecture-specific components. The same procedure is applied
unchanged to every model, providing a common measure of realised motion.
Because it requires only unlabelled video, it is also a supervision signal
available to other world-model systems rather than an advantage specific to
ours.

\paragraph{GT continuation can masquerade as control.}
Figure~\ref{fig:following_dir} compares recorded commands (GT) with
counterfactual commands (given by FLIP, these commands are generated by rotating the GT commands by 180 degrees in the action space so forward becomes backward and left becomes right) during training. GT agreement is
high almost immediately because the pretrained video model continues the
motion already present in the context. FLIP exposes whether the command
actually changes the continuation: steering is learned quickly, while genuine
forward--backward control emerges much later. This effect can be seen clearer when considering a variant that did not learn action conditioning but did learn to follow video as we show in Appendix~\ref{app:extended_control_following}. We see that the model first learns to complete videos, then it learns to be mean following when learning actions, and then it slowly learns the rarer actions, such as backwards. The
diagonal panels also align with their compound commands, showing that the
model composes throttle and steering rather than learning each direction as
an unrelated category.

\begin{figure}[t]
\centering
\includegraphics[
    width=0.4\textwidth
]{AnonymousSubmission/images/following_8node8pca_by_REALdir_select.png}
\caption{
    Training-time agreement between commanded and realized motion. N denotes how many minutes of each direction are present in the test set. GT uses the
    recorded command, while FLIP applies the opposite command to
    the same source context. Positive FLIP agreement indicates genuine command
    response rather than passive continuation. The full eight-direction plot
    appears in Fig.~\ref{fig:following_8pca}. Curves use a $9$-point rolling mean;
    $n$ denotes the number of evaluated chunks in each source-motion group.
}
\label{fig:following_dir}
\end{figure}

\paragraph{Response follows the data support.}
Figure~\ref{fig:response_flip} evaluates the settled response of the ablation
variants across command strength. Steering remains approximately
linear through the entire range, whereas throttle is less regular and saturates beyond the range
observed during training. As batch size or the number of critic-supervised PCA
dimensions is reduced, the reverse curves also bend back toward zero at
progressively larger negative command magnitudes. The ablations
therefore somewhat contract the model's range of control extrapolation. This highlights a limitation of our data-derived
approach: response strength depends strongly on the coverage of the observed
motion distribution. Conversely, the absence of the same contraction on the forward branch is consistent with its much denser representation in the training data.
Nevertheless, every functional variant recovers the sign of backward motion
despite its extreme scarcity; only the tokens-only model fails. AdaLN is also shown to be
essential for establishing control, while larger batches and richer critic
supervision improve how strongly and how far across the observed range that
control is maintained.

\begin{figure}[t]
    \centering
    \begin{subfigure}[t]{\columnwidth}
        \centering
        \includegraphics[
            width=0.7\linewidth
        ]{AnonymousSubmission/images/response_complete.png}
        \caption{Forward and backward response.}
        \label{fig:response_forward}
    \end{subfigure}

    \begin{subfigure}[t]{\columnwidth}
        \centering
        \includegraphics[
            width=0.7\linewidth
        ]{AnonymousSubmission/images/response_complete_steer.png}
        \caption{Steering response.}
        \label{fig:response_steer}
    \end{subfigure}

    \caption{
    Settled throttle and steering responses of the ablation variants, measured
    by the CoTracker--PCA readout. Most points use sign-flipped commands on
    held-out ride windows. Because genuine reverse windows are sparse, the
    positive-throttle sweep instead applies randomly sampled forward commands
    to held-out no-op contexts. Dashed vertical lines mark the approximate
    action range observed during training.
    }
    \label{fig:response_flip}
\end{figure}

\paragraph{Backward composition distinguishes the models.}
Figure~\ref{fig:wedge_all_g0} applies the same eight-direction action-forcing
test to all models. Many baselines fail to produce backward motion at all.
Others reverse under the pure backward command but lose or flip that component
when steering is added, showing that backward, backward-left, and
backward-right are not represented compositionally. Our model also reflects
the severe forward--backward data imbalance, producing asymmetric wedges,
but it gives the strongest and most consistent reverse response across both
pure and diagonal commands. minWM also preserves the correct structure, but
with substantially weaker realized motion.

\begin{figure}[t]
\centering
\includegraphics[
    width=0.47\textwidth
]{AnonymousSubmission/images/wedge_all_g0.png}
\caption{
    Realized throttle $g_0$ under eight commands. Wedges show $32$ scenes;
    blue denotes forward and red backward response. Outlines and dashed arcs
    show medians and interquartile ranges. Several models fail to reverse or
    lose reverse motion under diagonal commands. Independent radial scales
    make sign and shape, not magnitude, comparable.
}
\label{fig:wedge_all_g0}
\end{figure}

\subsection{Quality}

Action forced rollouts do not have paired ground-truth continuations, which makes standard reference-based video metrics unsuitable. We therefore evaluate quality through four complementary failure signals: style shift, implausibility, scene relocation, and geometric corruption. A model can avoid visible errors by remaining static. We therefore exclude near-static rollouts from all of the above, reusing the scene relocation metric. We detail the evaluation methods below but more information is provided in Appendix~\ref{app:quality_metrics}. 

Following prior VLM-based evaluation of generated visual content
\cite{videophy,vqa,videoscore,videobench}, we formulate each perceptual
criterion as a narrow binary visual question. Our use of answer token
probabilities follows VideoPhy and VQAScore~\cite{videophy,vqa}, avoiding
the ambiguity of parsing free-form explanations. We use Qwen3-VL-8B
\cite{QWEN}, conditioned on the last real frame of the input context and the
generated continuation. The reference and generated frames are explicitly
labelled, and the model assesses style shift, introduced objects, and severe
reality-breaking artifacts. We flag a probe when
$P(\mathrm{Yes}) > 0.5$.

Scene preservation is measured using ORB feature matches between the
generated end frame and either the real reference or sibling end frames.
RANSAC fits a projective image transformation to each pair and counts the
geometrically consistent matches. A rollout is marked as
relocated when fewer than $50$ verified inliers remain. The relocation
verifier agrees with human labels on $87.5\%$ of $80$ annotated rollouts,
and a further stratified visual audit obtains $89$--$96\%$ agreement across
the individual instruments. Using this metric, we define a rollout as static when its first and last generated frames retain more than $600$ RANSAC-verified feature matches. This filter removes $22\%$ of minWM
rollouts and at most $10\%$ for every other model.

Egomotion models must repeatedly re-render scene geometry as the viewpoint
changes. Their characteristic failure is therefore not always blur or global style drift, but localized over-warping: newly generated objects retain plausible colours, edges, and coarse shapes while their internal structure becomes melted, duplicated, or geometrically impossible. The other no reference methods find this class of error hard to qualify, so we use an alternative formulation here.

We instead combine three artifact-oriented signals. A
VLM~\cite{QWEN} judges whether structures are melted, warped, or
geometrically impossible, using answer-token probabilities rather than
generated text as before ($0.74$ AUC against a human evaluated subset). A generative artifact segmenter, PAL4VST~\cite{pal}, contributes the largest per-frame artifact fraction ($0.70$ AUC). Finally, we measure abnormal curvature in a monocular depth field produced by Depth Anything V2~\cite{depthanything} ($0.75$ AUC). After scene-relative standardization, their unweighted mean reaches $0.89$
AUC for structural corruption. The signals are therefore complementary: the VLM
and depth curvature respond strongly to impossible geometry, while artifact
segmentation detects localized synthetic defects.

All three signals are evaluated relationally. For each initial scene, a
rollout is compared with its sibling group, comprising the corresponding
generations from all models and ablations together with real FrodoBots
footage. Subtracting the sibling median reduces sensitivity to scene
difficulty and to naturally occurring defects in the real-world footage. We
flag a rollout when its scene-relative ensemble score exceeds the
$95$th percentile measured on real FrodoBots clips, fixing the reference
false-positive rate at approximately $5\%$ rather than selecting a threshold
from generated failures. No real clip approaches the scores of the
corrupted generations.

Table~\ref{tab:quality_all} shows that the models fail in distinct ways.
WorldPlay, Matrix-Game, and WorldCam frequently change visual style, produce
implausible content, relocate the scene, and corrupt geometry. Astra is less
extreme but remains weak across all four measures. Yume largely preserves
style and plausibility, yet relocates the scene in $60\%$ of active rollouts.

minWM is visually stable and usually preserves the scene, although $22\%$ of
its rollouts are excluded as near-static and its action response remains weak.
Our models remain active while achieving the lowest geometric-corruption rates
and low style and relocation error. The \texttt{pca8} model provides the
strongest overall balance, combining reliable backward and compositional
control with $3\%$ style shift, $16\%$ implausibility, $11\%$ relocation, and
$16\%$ geometric corruption.

Because our model is trained on FrodoBots footage while the baselines are not,
these results measure robustness on this real-world deployment distribution
rather than a domain-neutral ranking. They nevertheless demonstrate that a
rollout can remain sharp while changing style, leaving the scene, or
collapsing geometrically.

\begin{table}[t]
\centering
\small
\setlength{\tabcolsep}{4pt}
\caption{
    Generation-quality failure rates; lower is better. Style,
    Implausible, and Relocated are percentages of active rollouts from
    $240$ feature-valid examples per model. Geometry is the percentage of
    all $256$ rollouts exceeding the real-footage $95$th-percentile
    structural-corruption threshold.
}
\label{tab:quality_all}
\begin{tabular}{lrrrr}
\toprule
Model & Style\,\% & Implausible\,\% & Reloc. \,\% & Geometry\,\% \\
\midrule
WorldPlay     & 72 & 81 & 80 & 57 \\
Matrix-Game   & 77 & 72 & 95 & 51 \\
WorldCam      & 91 & 75 & 88 & 71 \\
Astra         & 29 & 54 & 66 & 45 \\
Yume          & \textbf{7} & \textbf{8} & 60 & 26 \\
minWM         & \textbf{1} & \textbf{2} & \textbf{5} & \textbf{18} \\
Ours (pca8)   & \textbf{3} & 16 & 11 & \textbf{16} \\
Ours (Batch64)& \textbf{7} & 24 & \textbf{8} & \textbf{16} \\
\bottomrule
\end{tabular}
\end{table}

Reduced PCA supervision increases implausibility and geometric corruption,
while smaller batches primarily increase relocation and high-frequency
degradation; full ablations are reported in
Appendix~\ref{app:quality_ablations}.

\paragraph{The zero command must stop the camera without freezing the world.}
A no-op command creates a distinct evaluation problem. The camera should
remain stationary, but pedestrians, traffic, vegetation, and other independent
scene dynamics should continue. Camera stability alone therefore rewards the
degenerate strategy of repeating the final context frame.

We evaluate no-op continuations against their paired real futures using
camera drift, residual scene motion, CD-FVD, and LPIPS. Camera drift measures
whether the null action is followed, while the residual-motion ratio compares
independent motion after removing the dominant background transformation.
We include a freeze baseline that repeats the final context frame, exposing
metrics that reward inactivity. All models and references are truncated to
the same temporal horizon.


\section{Conclusion \& Limitations}

We show that using a simple pixel tracker and finding the principal components of the pixel vectors results in recovering a reversible, composable, and linear action space with which we learn reliable backward control and composable steering with less than 1\% of our training set being backward data. Other models are shown to avoid rare egomotion such as backwards.

Our method is both strengthened and limited by it's core data dependency - any axes of egomotion present in the data can be recovered but motion not present in the data cannot be be learnt in this manner. We also find that while our model interpolates within the support provided, extrapolation performance is somewhat lacking.

Action forcing evaluation also shows that conventional evaluation of reference based quality is insufficient as a metric to evaluate world models. Our controllability wedges and quality metrics offer one potential manner in which WM evaluation can be done.

As a final takeaway, we wish to highlight that simpler methods that scale to the provided data are available and world model training no longer needs to be bounded by the quality of provided action labels, nor does it need to be relegated to those who can collect or label copious amounts of data.

\newpage
\bibliography{aaai2026}

\appendix

\section{Recovered Motion Basis}
\label{app:pca_components}

Figure~\ref{fig:pca_act} visualizes the leading motion directions recovered
by the fixed CoTracker--PCA pipeline. The first component, $g_0$, captures
longitudinal forward--backward motion; $g_1$ captures yaw; and $g_6$ captures
roll-like camera motion. Only $g_0$ and $g_1$ are exposed as control inputs.
The leading eight components are retained as critic targets so that secondary
camera and scene dynamics can regularize generation without enlarging the
two-dimensional control interface.

\begin{figure*}[t]
    \centering
    \includegraphics[
        width=0.95\textwidth
    ]{AnonymousSubmission/images/pca_components_flowfields.png}
    \caption{
        Principal motion components recovered from the fixed
        CoTracker--PCA pipeline. Each panel shows the displacement field
        associated with one PCA coordinate. The exposed control space uses
        $g_0$ for longitudinal motion and $g_1$ for yaw; $g_6$ captures
        roll-like motion. The remaining components describe secondary
        camera and scene dynamics and are used only for critic supervision.
    }
    \label{fig:pca_act}
\end{figure*}

\section{Cross-Model Evaluation Protocol}
\label{app:evaluation_protocol}

\subsection{Action-Forcing Benchmark}
\label{app:action_forcing_details}

\paragraph{Low-egomotion context selection.}
The benchmark uses $32$ contexts drawn from distinct held-out FrodoBots rides.
For each candidate window $w$, we compute the mean ego-motion magnitude across
its $N_w$ chunks:
\[
e(w)
=
\frac{1}{N_w}
\sum_{t=1}^{N_w}
\sqrt{
a_{t,\mathrm{throttle}}^2
+
a_{t,\mathrm{yaw}}^2
},
\]
where $a_{t,\mathrm{throttle}}$ and $a_{t,\mathrm{yaw}}$ are the
tanh-squashed PCA action coordinates for chunk $t$.

Let $q_{25}$ denote the $25$th percentile of the ego-motion scores across all
candidate windows. We define the eligibility threshold as
\[
\tau_{\mathrm{stat}}
=
\min(0.12,q_{25}),
\]
and retain candidate windows satisfying
\[
e(w)\leq\tau_{\mathrm{stat}}.
\]
This selects at most the lowest-motion quarter of the candidate population
while also imposing an absolute upper limit of $0.12$ in squashed PCA units.

The selection criterion restricts camera ego-motion rather than all visible
motion. A context may still contain independently moving pedestrians,
vehicles, vegetation, or other scene elements. A separate point-track
scene-activity score is used elsewhere for evaluation but is not part of
this low-egomotion selection criterion.

\paragraph{Commands.}
Each context is continued under the eight constant compass commands
\[
\{
\mathrm{F},\mathrm{FR},\mathrm{R},\mathrm{BR},
\mathrm{B},\mathrm{BL},\mathrm{L},\mathrm{FL}
\},
\]
giving $256$ rollouts per model. The same source contexts are used for every
model, and each command is held constant throughout the generated
continuation.

For our model, actions are ordered as
\[
a=(a_{\mathrm{throttle}},a_{\mathrm{yaw}}).
\]
All commands have Euclidean norm $s=0.5$. The cardinal commands are
\[
\mathrm{F}=(s,0),\qquad
\mathrm{B}=(-s,0),\qquad
\mathrm{R}=(0,s),\qquad
\mathrm{L}=(0,-s),
\]
and the diagonal commands are
\[
\begin{aligned}
\mathrm{FR}&=\frac{s}{\sqrt{2}}(s,s),&
\mathrm{FL}&=\frac{s}{\sqrt{2}}(s,-s),\\
\mathrm{BR}&=\frac{s}{\sqrt{2}}(-s,s),&
\mathrm{BL}&=\frac{s}{\sqrt{2}}(-s,-s).
\end{aligned}
\]
Thus, each active diagonal coordinate has magnitude
$s/\sqrt{2}\approx0.354$. Normalizing the diagonals in this way tests
translation--steering composition without increasing total command magnitude
relative to the cardinal directions.

\paragraph{Context and rollout duration.}
Our model receives one temporal context chunk containing three latent frames,
corresponding to $12$ RGB frames or $0.75$\,s at $16$\,fps. It then generates
eight autoregressive chunks, corresponding to $96$ RGB frames or $6.0$\,s.
The complete context-plus-generation clip therefore contains $108$ RGB frames
and spans $6.75$\,s.

The cross-model benchmark evaluates a common six-second generated horizon.
Each baseline is run through its native autoregressive interface and at its
native output frame rate.

\paragraph{Sampling and reproducibility.}
Our model uses the rectified flow-matching sampler with $48$ denoising steps
per generated chunk, \texttt{timestep\_shift}$=5.0$, and
\texttt{real\_guidance\_scale}$=0.0$, corresponding to no classifier-free
guidance. We evaluate the step-$5{,}000$ checkpoint
\texttt{causal\_lora\_step0005000.pt}. The same sampling configuration is used
for all of our ablations.

Sampling is deterministic. For generated block
$b\in\{0,\ldots,7\}$, the initial noise is drawn using
\[
\mathrm{seed}(b)=1234+b.
\]
The same fixed eight-block noise sequence is reused for every context and
command, holding sampling noise constant across the action-forcing
comparisons.

Each baseline uses its released evaluation checkpoint, sampler, inference-step
count, guidance configuration, and random-seed handling without modification.
The only extensions are the control-interface mappings described in
Section~\ref{app:baseline_controls}.

\paragraph{Evaluation resolution and motion readout.}
Each model is evaluated at its native output resolution; no shared resizing is
applied before the CoTracker--PCA controllability readout. The resolutions are
listed in Table~\ref{tab:action_forcing_resolutions}.

\begin{table}[t]
\centering
\small
\caption{Native output resolutions used by the action-forcing benchmark.}
\label{tab:action_forcing_resolutions}
\begin{tabular}{lc}
\toprule
Model & Output resolution \\
\midrule
Ours & $832\times480$ \\
minWM & $832\times480$ \\
WorldCam & $832\times480$ \\
HY-World 1.5 WorldPlay & $832\times480$ \\
Astra & $832\times480$ \\
Matrix-Game 2.0 & $640\times352$ \\
Yume-1.5 & $1280\times704$ \\
\bottomrule
\end{tabular}
\end{table}

The frozen CoTracker--PCA readout is applied independently at each model's
native resolution. Because the conditioning interfaces and native output
resolutions are not metrically matched, raw PCA response magnitudes are not
treated as directly comparable across models. Cross-model analysis instead
emphasizes response sign, requested direction, within-model
forward--backward asymmetry, and translation--steering composition. For the
separate cross-model image-translation analysis, only width-normalized
horizontal and vertical translations are compared.

\paragraph{Failed generations.}
Each context--command rollout is executed independently. If generation fails,
the rollout is skipped entirely: no video or metrics row is written, and the
missing result is not zero-filled or replaced. In practice, after the
control-interface implementations were finalized, all $256$ rollouts
completed successfully for every evaluated model. No benchmark rows were
therefore removed because of generation failure.


\subsection{Baseline Control Interfaces}
\label{app:baseline_controls}

\paragraph{Common protocol.}
Each compass command is translated into the closest operation supported by the
released model and repeated throughout generation. For pose-conditioned
models, the running camera pose is composed with a fixed relative translation
and/or yaw increment. For keyboard--mouse and text-conditioned models, the
corresponding control input or caption is held constant throughout the
rollout. For a given model, all parameters are fixed across scenes and
directions.

The interfaces do not share common physical units. A numerical translation or
rotation used by one model therefore cannot be interpreted as the same
physical command in another model. The benchmark evaluates whether each model
follows the direction and compositional structure of the command expressed
through its own released interface.

\paragraph{minWM.}
minWM accepts an autoregressive sequence of relative $\mathrm{SE}(3)$ camera
poses. We use the released trajectory-generator increments of $0.08$
translation and $3^\circ$ yaw per generation step. Forward and backward use
positive and negative longitudinal translation, left and right use signed
yaw, and diagonal commands combine the corresponding translation and yaw
increments.

\paragraph{HY-World 1.5 WorldPlay.}
WorldPlay is also driven using relative $\mathrm{SE}(3)$ camera poses and uses
the released trajectory scale of $0.08$ translation and $3^\circ$ yaw per
step. Its pose-string interface cannot express simultaneous translation and
yaw, so diagonal commands are supplied through the released JSON pose
interface. Backward and backward-diagonal commands use negative longitudinal
translation.

\paragraph{Astra.}
Astra conditions on a per-frame relative camera pose represented by a
flattened $3\times4$ matrix. We follow the scale used in the released
demonstration code, including its translation increment of $0.03$. The release
provides forward and signed-steering primitives. We construct backward motion
by reversing the longitudinal translation and construct backward-diagonal
commands by combining this translation with the released signed-yaw
primitive.

\paragraph{WorldCam.}
WorldCam is conditioned on per-pixel-frame camera-to-world trajectories in
OpenCV convention. Its pose scale is inherited from its training trajectories
and is not metric. Because the release does not include an eight-direction
trajectory generator, we construct one through the released pose interface.

The conditioning span uses identity poses. During generation, the running pose
is composed multiplicatively as
\[
C_{t+1}=C_tD,
\]
where $D$ contains a longitudinal translation of
\[
\Delta z=\pm0.02
\]
camera-to-world units per pixel frame and/or a yaw rotation of
$\pm1.5^\circ$ about the $y$-axis. Positive and negative translations are used
for the forward and backward command families, respectively. Positive yaw is
used for right-family commands and negative yaw for left-family commands.
Diagonal commands combine the corresponding translation and yaw within the
same increment.

The released runner performs $50$ autoregressive steps and produces $200$
generated frames. We use the fixed intrinsics
\[
f_x=f_y=723,\qquad c_x=960,\qquad c_y=540.
\]
The translation unit is arbitrary because WorldCam's pose scale derives from
its training trajectories. The translation and yaw increments were fixed
during interface calibration before the action-forcing benchmark and were
then applied unchanged across every scene and direction.

\paragraph{Matrix-Game 2.0.}
Matrix-Game accepts a six-dimensional binary keyboard stream and a
two-dimensional mouse stream, both held constant throughout generation. We
use its released benchmark mouse-yaw magnitude of $0.1$ per pixel frame. The
eight command mappings are given in
Table~\ref{tab:matrix_game_commands}.

\begin{table}[t]
\centering
\small
\caption{Matrix-Game 2.0 controls used for the eight-direction benchmark.}
\label{tab:matrix_game_commands}
\begin{tabular}{ccc}
\toprule
Command & Keyboard input & Mouse yaw \\
\midrule
$\mathrm{F}$  & \texttt{forward}=1  & $0$ \\
$\mathrm{B}$  & \texttt{backward}=1 & $0$ \\
$\mathrm{R}$  & none                & $+0.1$ \\
$\mathrm{L}$  & none                & $-0.1$ \\
$\mathrm{FR}$ & \texttt{forward}=1  & $+0.1$ \\
$\mathrm{FL}$ & \texttt{forward}=1  & $-0.1$ \\
$\mathrm{BR}$ & \texttt{backward}=1 & $+0.1$ \\
$\mathrm{BL}$ & \texttt{backward}=1 & $-0.1$ \\
\bottomrule
\end{tabular}
\end{table}

Keyboard controls are binary, whereas mouse yaw has a continuous magnitude.
Backward input is supported by the released interface but is not exercised in
the model's released benchmark protocol.

\paragraph{Yume-1.5.}
Yume is controlled through natural-language captions drawn from its supported
training vocabulary. Each command uses the common prefix
\begin{quote}
\small
\emph{This video depicts a city walk scene with a first-person view (FPV).}
\end{quote}
followed by the corresponding person and camera clauses in
Table~\ref{tab:yume_commands}. The complete caption is repeated throughout
the rollout. Yume exposes directional commands but no numerical command
magnitude and therefore cannot be evaluated using a continuous dose sweep.

\begin{table}[t]
\centering
\small
\caption{Yume-1.5 caption clauses used for the eight-direction benchmark.}
\label{tab:yume_commands}
\begin{tabular}{lll}
\toprule
Command & Person clause & Camera clause \\
\midrule
F  & moves forward (W)              & remains still \\
B  & moves backward (S)             & remains still \\
L  & stands still                    & turns left \\
R  & stands still                    & turns right \\
FL & moves forward and left (W+A)   & turns left \\
FR & moves forward and right (W+D)  & turns right \\
BL & moves backward and left (S+A)  & turns left \\
BR & moves backward and right (S+D) & turns right \\
\bottomrule
\end{tabular}
\end{table}

\paragraph{Authored interface extensions.}
Our additions operate only through the models' released conditioning
interfaces. We construct backward and backward-diagonal pose sequences for
Astra, backward--yaw combinations for Matrix-Game diagonals, an
eight-direction trajectory generator for WorldCam, and simultaneous
translation--yaw commands through WorldPlay's JSON interface. For the
stationary evaluation we further supply each model with a no-op command
through its native interface: an all-zero keyboard and mouse-control stream
for Matrix-Game, identity pose sequences for the camera-conditioned models
(WorldCam, WorldPlay, minWM), and a ``remains stationary'' prompt for Yume.
These inputs are representable by the released conditioning tensors but are
not part of the models' published benchmark protocols, which expose no named
or evaluated no-op command; we therefore report the resulting rollouts as
each model's response to an authored zero-input extension. No model
architecture, learned weight, or conditioning representation is modified.

\subsection{Response-Curve Protocol}
\label{app:response_details}

\paragraph{Ablation curves.}
The reverse-throttle and steering branches are measured using held-out ride
windows whose recorded commands are replaced by their rotated
counterfactuals. We use evaluations from the final $1{,}000$ training steps,
restrict the window offset to at most $81$, and measure only settled chunks by
discarding the first two generated chunks of each rollout. This is because there is a finite acceleration in the data and rapid direction changes induce a penalty from this acceleration. We find this in a lot of the other baseline models too, so all models use this protocol.

For each action axis, requested values are divided into $13$ linearly spaced
bins spanning the $2$nd to the $98$th percentile of the requested-command
distribution. A bin is plotted only when it contains at least $12$ samples.
Each plotted point is the mean realized CoTracker--PCA response within that
bin. Shaded regions show one standard error of the mean:
\[
\mathrm{SEM}
=
\frac{\operatorname{std}(x)}{\sqrt{n}}.
\]

Genuine reverse travel is too rare to produce a sufficiently populated
rotated forward branch. We therefore measure forward response using
held-out no-op contexts and the fixed positive throttle strengths
\[
\{0.1,0.2,\ldots,0.8\}.
\]
Settled response is again measured after discarding the first two generated
chunks. When more than $132$ samples are available at a command strength, we
subsample $132$ examples using a fixed random seed of $0$. This cap matches
the size of the largest FLIP bin.

The forward and reverse branches therefore probe the same throttle coordinate
but do not use identical starting-context distributions. They are used to
evaluate response sign, smoothness, saturation, and approximate gain rather
than as a paired forward--backward comparison. Dashed plot boundaries mark
the approximate support of the training actions, extending to approximately
$+0.5$ for forward throttle and $-0.3$ for reverse throttle.

\paragraph{Cross-model curves.}
Cross-model response curves use the same $32$ held-out contexts as the
eight-direction action-forcing benchmark. The compass commands produce the
five signed per-axis protocol levels
\[
\left\{
-0.5,\,
-\frac{0.5}{\sqrt{2}},\,
0,\,
\frac{0.5}{\sqrt{2}},\,
0.5
\right\}
\approx
\{-0.5,-0.354,0,0.354,0.5\}.
\]
For each level, the plotted point is the mean response over the $32$ scenes
and the shaded region shows $\pm1$ SEM.

Each baseline receives the corresponding fixed-strength command through its
native interface. The horizontal coordinate represents the signed protocol
level rather than a common physical control magnitude. The curves are
therefore interpreted through response sign and within-model
forward--backward asymmetry rather than absolute cross-model gain. The dotted
diagonal denotes unit gain.

\section{Extended Controllability Results}
\label{app:extended_control}

This section expands the controllability analysis from the main paper. We
first study how the conditioning pathway, batch size, and dimensionality of
critic supervision affect the acquisition of command control. We then report
the settled directional responses of the ablation variants, followed by
additional cross-model analyses of throttle, steering, and the roll-like
motion coupled to yaw.

\subsection{Conditioning-Pathway Ablation}
\label{app:conditioning_ablation}

Figure~\ref{fig:abl_injection} isolates the effect of the action-conditioning
pathway. Injecting actions through both AdaLN and action tokens produces the
fastest convergence and strongest counterfactual response. AdaLN alone still
establishes control, but converges more slowly and reaches a weaker response.
Action tokens alone appear competitive under GT evaluation but fail under
FLIP, showing that continuation of the source motion can be mistaken for
command following.

\begin{figure*}[t]
    \centering
    \includegraphics[
        width=0.82\textwidth
    ]{AnonymousSubmission/images/following_FAMILY_injection_by_REALdir.png}
    \caption{
        Conditioning-pathway ablation under counterfactual FLIP evaluation,
        grouped by source-ride motion direction. Dual injection through AdaLN
        and action tokens converges fastest and reaches the strongest response.
        AdaLN-only conditioning learns control more slowly, while tokens-only
        conditioning remains anti-aligned on throttle and unstable on steering.
        Panels and smoothing follow Fig.~\ref{fig:following_dir}.
    }
    \label{fig:abl_injection}
\end{figure*}


\subsection{Training-Time Command Following}
\label{app:extended_control_following}

\paragraph{Default model.}
Figure~\ref{fig:following_8pca} extends the main-paper training analysis to all
eight source-motion groups for the default pca8 model with global batch size
$32$. GT agreement rises early because the pretrained model can continue the
motion already present in the context. FLIP agreement instead measures whether
the supplied action changes that continuation. Steering is acquired first,
while backward and backward-diagonal control emerge later.

\begin{figure*}[t]
    \centering
    \includegraphics[
        width=0.82\textwidth
    ]{AnonymousSubmission/images/following_8node8pca_by_REALdir.png}
    \caption{
        Training-time command following for the default pca8 model, measured as
        cosine agreement between the commanded action and realized
        CoTracker--PCA motion. Panels group generated chunks by the motion
        direction of the source ride. Blue uses the recorded command (GT);
        red applies its sign-flipped counterfactual (FLIP) to the same context.
        GT agreement can be high under ordinary video continuation, whereas
        FLIP reveals when the generated motion becomes responsive to the
        command. Backward control emerges last, while diagonal agreement
        demonstrates composition of throttle and steering. Curves use a
        $9$-point ($\approx450$-step) centered rolling mean.
    }
    \label{fig:following_8pca}
\end{figure*}

\paragraph{Tokens-only conditioning.}
Figure~\ref{fig:noadaln_gt_flip} illustrates why GT agreement is insufficient
on its own. The no-AdaLN, tokens-only variant obtains high GT agreement in
several common source-motion contexts, but its FLIP response remains negative
or unstructured. It therefore continues motion inherited from the context
without acquiring reliable command control.

\begin{figure*}[t]
    \centering
    \includegraphics[
        width=0.82\textwidth
    ]{AnonymousSubmission/images/following_noadaln_by_REALdir.png}
    \caption{
        GT and FLIP agreement for the no-AdaLN, tokens-only variant. GT
        agreement is high in common source-motion contexts, but FLIP remains
        negative or unstructured. The variant therefore continues motion
        already present in the context without acquiring reliable command
        control. Panel titles denote the motion direction of the source
        context rather than the supplied counterfactual command.
    }
    \label{fig:noadaln_gt_flip}
\end{figure*}

\paragraph{Larger-batch model.}
Figure~\ref{fig:following_16pca} shows the corresponding training curves for
the Batch64 model. GT agreement again rises rapidly, but FLIP agreement also
becomes consistently positive across all eight directions. Relative to the
default batch size, the larger batch produces smoother training dynamics and
a stronger response in the harder throttle directions.

\begin{figure*}[t]
    \centering
    \includegraphics[
        width=0.82\textwidth
    ]{AnonymousSubmission/images/following_16node_by_REALdir.png}
    \caption{
        GT and FLIP agreement for the Batch64 model. Although GT agreement
        rises rapidly by continuing motion already present in the source
        context, FLIP agreement also becomes strongly positive across all
        eight directions, demonstrating genuine command response. Panel titles
        denote the motion direction of the source context rather than the
        supplied counterfactual command.
    }
    \label{fig:following_16pca}
\end{figure*}


\subsection{Batch-Size and Critic-Supervision Ablations}
\label{app:control_training_ablations}

Reducing either global batch size
(Fig.~\ref{fig:abl_nodes}) or the number of PCA components supervised by the
critic (Fig.~\ref{fig:abl_encoders}) delays convergence and weakens the
realized reverse response, but does not prevent recovery of the backward
direction. Every functional conditioning variant recovers backward control by
step $5{,}000$; only the tokens-only model fails. The corresponding effects on
generation quality are reported in
Appendix~\ref{app:quality_ablations}.

\begin{figure*}[t]
    \centering
    \includegraphics[
        width=0.82\textwidth
    ]{AnonymousSubmission/images/following_FAMILY_nodes_by_REALdir.png}
    \caption{
        Batch-size ablation using global batch sizes $16$, $32$, and $64$
        under FLIP evaluation. All batch sizes recover the same directional
        structure by step $5{,}000$, but larger batches acquire the hardest
        throttle responses earlier and retain stronger reverse motion.
    }
    \label{fig:abl_nodes}
\end{figure*}

\begin{figure*}[t]
    \centering
    \includegraphics[
        width=0.82\textwidth
    ]{AnonymousSubmission/images/following_FAMILY_encoders_by_REALdir.png}
    \caption{
        Critic-supervision ablation using the leading $2$, $4$, or $8$ PCA
        components. Two supervised components suffice to recover directional
        control, while richer supervision improves convergence and realized
        response strength.
    }
    \label{fig:abl_encoders}
\end{figure*}


\subsection{Settled Ablation Responses}
\label{app:settled_ablation_responses}

Figures~\ref{fig:wedge_g0} and~\ref{fig:wedge_g1} compare the settled
directional responses of the ablation variants. Batch size has its largest
effect on reverse throttle: Batch16 is substantially weaker under backward and
backward-diagonal commands than Batch32 and Batch64. Steering is less
sensitive to batch size and critic-supervision dimensionality, with every
functional variant retaining the expected signed four-lobe structure. The
tokens-only model fails to organize either action axis.

\begin{figure*}[t]
    \centering
    \includegraphics[
        width=0.82\textwidth
    ]{AnonymousSubmission/images/wedge_g0.png}
    \caption{
        Settled throttle response $g_0$ by commanded direction for the
        ablation variants. Each translucent wedge represents one of $32$
        scenes; blue denotes forward-signed response and red denotes
        backward-signed response. Solid outlines show medians, and panel rims
        show the local $95$th percentile. Every functional variant recovers
        reverse motion, whereas the tokens-only model renders all commands as
        forward.
    }
    \label{fig:wedge_g0}
\end{figure*}

\begin{figure*}[t]
    \centering
    \includegraphics[
        width=0.82\textwidth
    ]{AnonymousSubmission/images/wedge_g1.png}
    \caption{
        Settled steering response $g_1$ by commanded direction for the
        ablation variants, rendered as in Fig.~\ref{fig:wedge_g0}. Every
        functional variant shows the expected signed four-lobe structure.
        Left-family responses are weaker than right-family responses,
        consistent with an asymmetry in the underlying tracked-motion data.
    }
    \label{fig:wedge_g1}
\end{figure*}


\subsection{Cross-Model Response Curves}
\label{app:cross_model_response}

Figure~\ref{fig:response_gain_all} provides a scalar view of the behavior
summarized by the main-paper directional wedges. The curve construction and
uncertainty estimates are defined in
Section~\ref{app:response_details}.

Our model supports continuous-valued actions, whereas several baselines expose
only fixed pose increments, binary controls, or text commands without a
numerical strength parameter; see
Section~\ref{app:baseline_controls},
Table~\ref{tab:matrix_game_commands}, and
Table~\ref{tab:yume_commands}. Baselines are therefore shown only at strengths
supported by their released interfaces, and unsupported intermediate values
are not interpolated.

Because the interfaces use different and often non-metric units, the curves
are interpreted primarily through response sign and within-model
forward--backward asymmetry rather than absolute cross-model gain.

\begin{figure}[t]
    \centering

    \begin{subfigure}[t]{\columnwidth}
        \centering
        \includegraphics[
            width=0.82\linewidth
        ]{AnonymousSubmission/images/response_curves_eval_throttle.png}
        \caption{Throttle response.}
        \label{fig:response_eval_forward}
    \end{subfigure}

    \begin{subfigure}[t]{\columnwidth}
        \centering
        \includegraphics[
            width=0.82\linewidth
        ]{AnonymousSubmission/images/response_curves_eval_steer.png}
        \caption{Steering response.}
        \label{fig:response_eval_steer}
    \end{subfigure}

    \caption{
        Cross-model action-forcing response curves from $32$ low-egomotion
        held-out contexts. Realized motion is measured using the common
        CoTracker--PCA readout, and shaded bands show $\pm1$ SEM. Baselines are
        shown only at strengths supported by their released interfaces.
        The dotted diagonal denotes unit gain as a visual reference.
    }
    \label{fig:response_gain_all}
\end{figure}

\subsection{Cross-Model Steering and Yaw--Roll Coupling}
\label{app:cross_model_steering_roll}

Figure~\ref{fig:wedge_all_g1} shows that every evaluated model recovers the
requested steering sign, although response magnitude and consistency across
scenes differ. Figure~\ref{fig:wedge_all_g6} measures the roll-like component
$g_6$, which is not directly commanded but responds systematically to yaw
across several models.

A non-zero $g_6$ response is not necessarily a failure. Real egocentric turns
are rarely pure yaw rotations: vehicle dynamics, suspension, camera mounting,
and body motion can introduce a correlated roll component. Because the PCA
basis is recovered from real video, $g_6$ may capture this naturally coupled
motion. We therefore use it as a diagnostic of yaw--roll coupling rather than
assuming that all roll should vanish. Large or inconsistent responses remain
undesirable, particularly when they exceed the coupling observed in real
motion.

\begin{figure*}[t]
    \centering
    \includegraphics[
        width=0.82\textwidth
    ]{AnonymousSubmission/images/wedge_all_g1.png}
    \caption{
        Steering response $g_1$ by commanded direction for all models.
        Rendering follows the main-paper throttle wedges. Every model recovers
        the expected steering sign, while response magnitude and variance
        differ across systems.
    }
    \label{fig:wedge_all_g1}
\end{figure*}

\begin{figure*}[t]
    \centering
    \includegraphics[
        width=0.82\textwidth
    ]{AnonymousSubmission/images/wedge_all_g6.png}
    \caption{
        Uncommanded roll response $g_6$ by commanded direction for all models.
        Smaller responses are preferable. Direction-dependent roll indicates
        coupling between the requested planar action and out-of-plane camera
        motion.
    }
    \label{fig:wedge_all_g6}
\end{figure*}
\section{Reference-Free Quality Evaluation Methodology}
\label{app:quality_metrics}

Action-forced rollouts do not admit a unique ground-truth continuation, making
paired reference metrics unsuitable. We instead evaluate complementary
failure modes that can be measured without a target video: visual style shift,
implausible content, scene relocation, geometric corruption, and progressive
loss of high-frequency detail. This section defines the evaluation population,
inputs, scoring rules, thresholds, and validation procedure for each
instrument.


\subsection{Evaluation Population}
\label{app:quality_population}

The action-forcing benchmark initially contains $240$ rollouts per model. A model may avoid visible generation failures by producing little or no
motion. We therefore exclude near-static rollouts. Static detection uses the
ORB matching and homography-verification pipeline defined in
Section~\ref{app:scene_relocation}. A rollout is classified as static when
its first generated frame and six-second-horizon frame retain more than
$600$ RANSAC-verified feature matches.

Style-shift, implausibility, relocation, and geometric-integrity rates are
reported over the remaining non-static rollouts, referred to as the active
evaluation population. The same canonical per-rollout active mask is used for
all four measures. The resulting denominators are reported in
Table~\ref{tab:quality_denominators}.

\begin{table}[t]
\centering
\small
\caption{
    Quality-evaluation populations. Each model contributes $240$
    feature-valid rollouts; Active additionally excludes near-static
    continuations using the ORB-based criterion.
}
\label{tab:quality_denominators}
\begin{tabular}{lrr}
\toprule
Model or variant & Total & Active \\
\midrule
WorldPlay          & 240 & 238 \\
Matrix-Game        & 240 & 240 \\
WorldCam           & 240 & 216 \\
Astra              & 240 & 215 \\
Yume               & 240 & 237 \\
minWM              & 240 & 187 \\
\midrule
Ours (pca8)        & 240 & 239 \\
Ours (pca4)        & 240 & 237 \\
Ours (pca2)        & 240 & 232 \\
Ours (16node)      & 240 & 237 \\
Ours (4node)       & 240 & 236 \\
Ours (noatok)      & 240 & 229 \\
Ours (noadaln)     & 240 & 200 \\
\bottomrule
\end{tabular}
\end{table}


\subsection{Style and Implausibility Probes}
\label{app:vlm_quality_probes}

We evaluate visual style shift and implausible content using two binary probes
implemented with Qwen3-VL-8B~\cite{QWEN}. Each probe receives one real
reference frame followed by $16$ generated frames distributed evenly over
the available continuation, allowing it to assess changes across the rollout
rather than at a single time point.

The reference is frame $8$ of the shared real context. Generated-frame
selection is deterministic: the first and final available frames within a
six-second window are included, with the remaining frames placed uniformly
between them using \verb|np.linspace(...).round().astype(int)|. The first
generated-frame indices are $9$ for ours, $4$ for Astra, $1$ for Matrix-Game,
$13$ for minWM, $65$ for WorldCam, and $1$ for WorldPlay and Yume. For our
$108$-frame rollouts, the selected indices are
$\{9,15,22,28,35,41,47,54,60,67,73,79,86,92,99,105\}$. Shorter baseline
rollouts are sampled over their available duration. All images are resized to
$640\times352$ and visibly labelled as \textsc{Reference} or
\textsc{Generated}.

Each probe begins with the following preamble:

\begin{quote}
\small
\emph{The first image labelled REFERENCE is the last real frame of a driving
video---the true scene. The remaining images labelled GENERATED are an AI
world model's continuation of that exact scene, in temporal order. The model
was supposed to continue the SAME scene in the SAME visual style with
plausible content.}
\end{quote}

The preamble is followed by one of the two questions below.

\paragraph{Style shift.}
\begin{quote}
\small
\emph{Does the visual STYLE of the generated frames depart from the
reference---e.g.\ becoming painted, game-like, cartoonish, watercolour,
oversaturated, or otherwise a different rendering style? Ordinary lighting
or exposure changes do NOT count. Answer Yes or No.}
\end{quote}

\paragraph{Implausible content.}
\begin{quote}
\small
\emph{Is there a SIGNIFICANT uncanny or reality-breaking failure in the
generated frames---impossible geometry, surfaces dissolving into abstract
patterns, large corrupted regions---that a casual viewer would notice within
one second? Ignore small local artifacts and minor blur. Answer Yes or No.}
\end{quote}

Each question ends with the literal instruction
\texttt{Answer with ONLY \{"answer": "Yes"\} or \{"answer": "No"\}.}
The assistant response is force-prefixed with
\verb|{"answer": "|. No free-form text is decoded. The probe score is the
softmax probability of \emph{Yes} against \emph{No} at the next-token
position, and a rollout is flagged when
\[
P(\mathrm{Yes})>0.5.
\]


\subsection{Scene-Relocation Instrument}
\label{app:scene_relocation}

Scene identity is measured independently of the VLM probes. Frames are
resized to $640\times352$, converted to grayscale, and processed using
\[
\texttt{cv2.ORB\_create(3000)}.
\]
Descriptors are matched using
\[
\texttt{BFMatcher(NORM\_HAMMING, crossCheck=True)}.
\]
Matching therefore uses brute-force Hamming distance with mutual cross-check;
no Lowe ratio test is applied.

Matches are geometrically verified using a single homography:
\[
\texttt{cv2.findHomography(src,dst,cv2.RANSAC,5.0)},
\]
where $5.0$ is the reprojection threshold in pixels. The verified-inlier
count is the sum of the returned RANSAC mask.

For each source scene, one joint sibling set is constructed from the
six-second-horizon end frame of every evaluated rollout: six external
baselines and seven ablation variants. The real reference is frame index $8$
from the shared source context.

For each rollout, we retain the largest number of RANSAC-verified ORB matches
obtained against either the real reference frame or the endpoint of any other
model generated from the same source scene. A rollout is classified as
relocated when this best match contains fewer than $50$ verified
correspondences.

Using the best match makes the test intentionally conservative: a continuation
is considered to remain in-scene whenever it agrees with at least one valid
view of the source scene, rather than requiring agreement with every sibling
rollout.

The static filter uses the same ORB, descriptor-matching, and homography
pipeline. It compares each rollout's first generated frame with its
six-second-horizon end frame and classifies the rollout as static when the
verified-inlier count exceeds $600$.


\subsection{Geometric-Integrity Instrument}
\label{app:geometry_instrument}

Egomotion models repeatedly re-render scene geometry as the viewpoint changes.
Their characteristic failure is often localized over-warping rather than
global blur: objects may retain plausible colour, texture, and coarse shape
while their internal structure becomes melted, duplicated, or geometrically
impossible. Whole-frame image statistics therefore often miss defects that
are immediately visible to humans.

We combine three artifact-oriented components: a VLM judgement of melted or
impossible structure, PAL4VST~\cite{pal} artifact area, and abnormal curvature
in a Depth-Anything-V2-Small~\cite{depthanything} relative-depth field.


\paragraph{VLM melt score.}

The deployed melt component uses Qwen2.5-VL-7B with the following prompt:

\begin{quote}
\small
\emph{Are any buildings, walls or structures in this image melted, warped,
smeared or geometrically impossible, like a corrupted AI-generated image?
Answer only Yes or No.}
\end{quote}

Six generated frames at indices
\[
\{24,40,56,72,88,104\}
\]
are evaluated. For each frame, the upper $60\%$ is retained and divided into
left and right halves, producing $12$ crops per rollout.

No text is decoded. At the final-token position, positive and negative token
logits are pooled separately:
\[
s_{\mathrm{yes}}
=
\operatorname{logsumexp}_{v\in\mathcal Y}\ell_v,
\qquad
s_{\mathrm{no}}
=
\operatorname{logsumexp}_{v\in\mathcal N}\ell_v,
\]
where
\[
\mathcal Y
=
\{\texttt{Yes},\texttt{\char32 Yes},\texttt{yes}\},
\qquad
\mathcal N
=
\{\texttt{No},\texttt{\char32 No},\texttt{no}\}.
\]
The crop score is
\[
P(\mathrm{Yes})
=
\sigma\!\left(s_{\mathrm{yes}}-s_{\mathrm{no}}\right),
\]
and the rollout-level component is the mean of the four largest crop scores.

% TODO (RUNNING):
% Insert the matched comparison of Qwen2.5-VL-7B, InternVL3-8B, and
% Cosmos-Reason1-7B once the benchmark completes.
%
% Report all three models under the identical prompt, frame selection, crop
% construction, token-logit scoring, and top-four aggregation rule. Include
% the exact checkpoint, positive and negative sample counts, development-set
% AUC, and any degenerate score behaviour. State that the prompt and
% aggregation were not tuned separately for each model.


\paragraph{PAL4VST artifact area.}

Eight wall-clock timestamps are selected uniformly over the generated span.
For our $16$\,fps rollouts, these correspond to frame indices
\[
\{26,37,49,60,72,84,95,107\}.
\]
For each frame, the upper $15\%$ is removed to exclude the distant haze band
and the remaining image is divided into left and right halves. Each half is
resized to $512\times512$ and processed by PAL4VST. The artifact-pixel
fraction is computed independently for each half, and the per-frame component
is their mean. The rollout-level component,
\texttt{pal4vst\_max}, is the maximum per-frame fraction over the eight
selected frames.


\paragraph{Depth-field curvature.}

Depth-Anything-V2-Small predicts a relative-depth field for three frames
selected uniformly from the base window near the generation hand-off. Each
prediction is resized to $208\times112$ and normalized by its own median:
\[
\widetilde d
=
\frac{d}{\operatorname{median}(d)+10^{-6}}.
\]
The frame-level curvature score is the mean absolute Laplacian:
\[
R(d)
=
\operatorname{mean}
\left|
\operatorname{Laplacian}(\widetilde d)
\right|.
\]
The rollout-level component, \texttt{depth\_rough\_base}, is the mean of
$R(d)$ over the three selected frames.

\paragraph{Sibling normalization and ensemble.}

The three components operate on different numerical scales and are strongly
affected by scene difficulty. We therefore normalize them in two stages.
First, each component is expressed relative to the median score of all
baseline and ablation rollouts generated from the same source scene. This
removes scene-specific effects, such as unusual architecture, camera noise,
or difficult lighting. Real footage is not included in this sibling median.

Second, each sibling-relative component is converted to a robust standardized
score using its fleet-wide median and median absolute deviation (MAD). This
places the VLM, PAL4VST, and depth signals on comparable scales while limiting
the influence of extreme failures. Concretely,
\[
z_{i,k}
=
\frac{
r_{i,k}-\operatorname{median}(r_{\cdot,k})
}{
1.4826\,\operatorname{MAD}(r_{\cdot,k})
},
\]
where \(r_{i,k}\) is the same-scene-relative value for rollout \(i\) and
component \(k\). The factor \(1.4826\) is the conventional scaling that makes
MAD comparable to standard deviation; the implementation adds a negligible
constant to the denominator only to prevent division by zero.

The final geometric-integrity score is the unweighted mean of the three
standardized components:
\[
G_i
=
\frac{
z_{i,\mathrm{VLM}}
+
z_{i,\mathrm{PAL}}
+
z_{i,\mathrm{depth}}
}{3}.
\]
Using equal weights avoids fitting the ensemble to the relatively small
labelled set.

We set the decision threshold to the \(95\)th percentile of scores measured
on real FrodoBots clips. This corresponds to \(G_i=0.826\) in the evaluated
fleet and fixes the reference false-positive rate at approximately \(5\%\).
Real footage is used only for this calibration and does not enter the sibling
normalization or reported denominators. Geometry rates use the same active
rollout population as the style, implausibility, and relocation measures.

\begin{table}[t]
\centering
\small
\caption{
    Candidate measures for severe structural corruption. AUCs are measured on
    the human-labelled ablation set.
}
\label{tab:geometry_metric_auc}
\begin{tabular}{lc}
\toprule
Signal & AUC \\
\midrule
Whole-frame and pixel measures       & $0.40$--$0.62$ \\
Line-segment collinearity            & $\leq0.65$ \\
PAL4VST artifact fraction            & $0.70$ \\
VLM melt score                       & $0.74$ \\
Depth-field curvature                & $0.75$ \\
Unweighted three-component ensemble  & $\mathbf{0.89}$ \\
\bottomrule
\end{tabular}
\end{table}

During validation, sibling medians were computed over the ablation variants
only; fleet deployment instead uses the joint baseline-and-ablation sibling
group. The labelled set was used to assess both the individual components and
their unweighted ensemble, so the reported AUC is a development-set estimate
rather than a nested held-out result.

The instrument targets severe structural collapse rather than subtle
inconsistencies. Large blobs, line tangles, and irregular warping are detected
reliably, while duplicated wires and small local surface defects remain
difficult.

\subsection{High-Frequency Degradation Instrument}
\label{app:high_frequency_metric}

Some variants progressively lose contrast and fine image structure during
long rollouts. We measure this effect using the change in Laplacian variance,
a standard indicator of high-frequency image content.

For a temporal window \(W\), let
\[
L_r(W)
=
\frac{1}{|W|}
\sum_{t\in W}
\operatorname{Var}\!\left(\nabla^2 g_{r,t}\right),
\]
where \(g_{r,t}\) is the grayscale frame from rollout \(r\) at time \(t\).
The base window contains four frames beginning near \(t=1\)\,s, while the
end window contains eight frames sampled with stride two near the end of the
rollout.

The loss of high-frequency detail in rollout \(r\) is
\[
\Delta_r
=
L_r(W_{\mathrm{base}})
-
L_r(W_{\mathrm{end}}),
\]
so that larger positive values indicate a greater reduction in sharpness.

To account for scene difficulty and source-image quality, we express this
loss relative to the other models generated from the same source scene:
\[
\mathcal B_r
=
\Delta_r
-
\operatorname{median}_{s\in\mathcal S_r}\Delta_s,
\]
where \(\mathcal S_r\) is the sibling group for rollout \(r\). Positive
\(\mathcal B_r\) therefore indicates greater sharpness loss than the typical
continuation of the same scene, while negative values indicate better
retention of high-frequency detail.

Against \(84\) human-labelled rollouts from \(12\) scenes, the continuous
score achieves AUC \(0.86\), compared with \(0.83\) for dark-channel rise and
\(0.79\) for contrast loss. To summarize the continuous score as a failure rate, we classify a rollout as
high-frequency degraded when $\mathcal B_r>150$. This fixed reporting
threshold was not optimized on the labelled evaluation set. The resulting
percentage is therefore a descriptive summary of the score distribution,
while the AUC evaluates the continuous metric independently of any particular
threshold. Dark-channel and contrast measurements are reported only as
alternative candidate metrics and do not contribute to this classification.

\paragraph{Dirty-source exclusion.}

A rollout is marked as dirty when its base-window Laplacian variance falls
below the fleet-wide $12$th percentile:
\[
L_r(W_{\mathrm{base}})
<
q_{12}
\left(
\{L(W_{\mathrm{base}})\}_{\mathrm{fleet}}
\right).
\]
Sibling-relative normalization is computed before this exclusion, so dirty
rollouts still contribute to the same-scene median. They are removed only
when aggregating the reported high-frequency degradation rates.


\subsection{Instrument Validation}
\label{app:quality_validation}

The scene-relocation instrument agrees with human labels on $87.5\%$ of
$80$ annotated rollouts and remains stable for thresholds between $43$ and
$60$ verified inliers. In a blind check on nine previously unseen videos,
eight predictions recorded before inspection were correct. A further
stratified audit of $28$ videos produced $89$--$96\%$ agreement across the
individual quality instruments, with disagreements concentrated in
borderline cases.


\subsection{Qualitative Instrument Audit}
\label{app:quality_qualitative_audit}

Figure~\ref{fig:quality_metric_audit} qualitatively audits both successful
detections and representative failure cases. Its purpose is to make the
operating range of each instrument visible and to show where its decisions
agree or disagree with human judgement.

\begin{figure*}[t]
\centering
\fbox{
\parbox[c][0.29\textheight][c]{0.96\textwidth}{
\centering
\textbf{Placeholder: qualitative quality-metric audit}\\[0.5em]
Real context and representative generated frames for clean, style-shift,
implausible, relocated, geometrically corrupted, high-frequency-degraded,
and detector-missed examples.
}}
\caption{
    Qualitative audit of the reference-free quality instruments. Each row
    shows real context and representative generated frames together with the
    human annotation, automatic decision, and relevant score. The examples
    include correctly accepted continuations, correctly detected failures,
    and representative borderline or detector-missed cases, illustrating
    both the intended operating range and the limitations of the instruments.
}
\label{fig:quality_metric_audit}
\end{figure*}

% TODO:
% Replace the placeholder with the selected human-audited examples.
%
% Include:
% 1. A clean true negative.
% 2. A correctly detected style shift.
% 3. A correctly detected implausible-content failure.
% 4. A correctly detected scene relocation, with consensus inliers shown.
% 5. A correctly detected geometric collapse, with the VLM, PAL4VST,
%    depth-curvature, and ensemble scores shown.
% 6. A correctly detected high-frequency degradation case, with
%    \mathcal B shown.
% 7. A representative false negative or borderline miss.
% 8. A false positive or threshold-borderline example, if available.
%
% For each example, show the real reference and three generated frames from
% the early, middle, and late portions of the continuation. Report the model,
% scene identifier, human label, metric decision, raw score, and threshold.
\section{Quality Evaluation Results}
\label{app:quality_results}

We apply the reference-free instruments defined in
Section~\ref{app:quality_metrics} to the external baselines and our ablation
variants. We first compare quality failures across models and then examine how
the action-conditioning design, critic supervision, and training batch size
affect the failure profile of our model.


\subsection{Cross-Model Quality}
\label{app:cross_model_quality}

The cross-model comparison measures style shift, implausible content, scene
relocation, and geometric corruption over the action-forcing evaluation set.
Style, implausibility, and relocation rates are computed over the active
populations reported in Table~\ref{tab:quality_denominators}.

% TODO (RUNNING):
% Insert the corrected cross-model table and associated plot after the
% reference-frame comparison has completed.
%
% The final table should report, for every external baseline and Ours (pca8):
%     Style (\%)
%     Implausible (\%)
%     Relocated (\%)
%     Geometric-integrity flag (\%)
%
% Update the table using the selected reference policy and state whether the
% corrected context boundary changes any model ranking.
%
% Suggested table label:
% \label{tab:quality_cross_model}
%
% Suggested plot label:
% \label{fig:quality_cross_model}

% TODO:
% Write the cross-model interpretation only after the corrected rates are
% available. Focus on differences in failure profile rather than reducing the
% comparison to a single aggregate quality ranking.

\subsection{Quality Ablations}
\label{app:quality_ablations}

Table~\ref{tab:quality_ablation} compares the failure profiles of the
ablation variants across style shift, implausible content, scene relocation,
and geometric corruption. The default \texttt{pca8} model provides the
strongest overall balance: it has the lowest style-shift rate among the
responsive variants, a low implausibility rate, and relatively low relocation
and geometric-corruption rates. The larger \texttt{16node} configuration
further reduces relocation, but does not improve the other quality measures.

Reducing the PCA supervision dimension primarily affects the realism and
structure of newly generated content. Relative to \texttt{pca8}, the
implausibility rate rises from $16\%$ to $35\%$ for \texttt{pca4} and
$39\%$ for \texttt{pca2}. The corresponding geometric-corruption rate rises
from $16\%$ to $28\%$ and $30\%$, respectively. This indicates that weaker
critic supervision permits continuations that remain visually coherent at a
coarse level while developing increasingly implausible or malformed scene
structure.

Reducing the training scale from \texttt{16node} to \texttt{4node} produces
a different failure profile. Relocation increases from $8\%$ to $24\%$, while
geometric corruption increases from $16\%$ to $23\%$. Removing the action
tokens produces similar rates of implausibility and geometric corruption, but
a lower relocation rate than \texttt{4node}.

The apparently low implausibility and geometry rates of \texttt{noadaln}
should not be interpreted as strong generation quality. As shown by the
controllability results, this variant frequently produces weak or
unresponsive continuations and therefore introduces less new scene structure
that can become visibly corrupted. It nevertheless exhibits the highest
relocation rate in the ablation grid.

\begin{table}[t]
\centering
\small
\setlength{\tabcolsep}{4pt}
\caption{
    Quality failure rates across the ablation grid; lower is better.
    Style, Implausible, and Relocated are percentages of active rollouts from
    the $240$ feature-valid examples per variant. Geometry reports the
    percentage flagged for severe structural corruption by the ensemble
    defined in Section~\ref{app:geometry_instrument}. Relocated denotes fewer
    than $50$ RANSAC-verified matches to the real reference or any sibling end
    frame. Geometry reports severe structural-corruption flags from the ensemble
    defined in Section~\ref{app:geometry_instrument}.
}
\label{tab:quality_ablation}
\begin{tabular}{lcccc}
\toprule
Variant
& Style\,\%
& Implausible\,\%
& Reloc.\,\%
& Geometry\,\% \\
\midrule
16node  & $7$ & $24$ & $8$  & $17$ \\
pca8    & $3$ & $16$ & $11$ & $16$ \\
pca4    & $9$ & $35$ & $14$ & $28$ \\
pca2    & $9$ & $39$ & $15$ & $32$ \\
4node   & $7$ & $29$ & $24$ & $24$ \\
noatok  & $7$ & $27$ & $14$ & $24$ \\
noadaln & $6$ & $12$ & $28$ & $19$ \\
\bottomrule
\end{tabular}
\end{table}


% TODO (RUNNING):
% Refresh the Geometry column if the matched VLM comparison changes the
% deployed geometry component or any ensemble classifications.
%
% Confirm and state the exact denominator used for the Geometry column.


\subsection{High-Frequency Degradation}
\label{app:high_frequency_results}

Table~\ref{tab:hf_losses} reports the sibling-relative high-frequency
degradation score $\mathcal B$ defined in
Section~\ref{app:high_frequency_metric}. Dirty-source rollouts are excluded
from these summaries.

The \texttt{16node}, \texttt{pca8}, and \texttt{pca4} variants remain
stable, exceeding the descriptive threshold $\mathcal B>150$ in only $2\%$
of evaluated rollouts. Reducing the PCA supervision dimension further to
\texttt{pca2} raises this rate to $20\%$, while removing the action tokens
raises it to $17\%$.

The strongest degradation occurs under the smaller \texttt{4node} training
configuration and in \texttt{noadaln}. These variants exceed the threshold
in $47\%$ and $59\%$ of rollouts, respectively. Their positive median
$\mathcal B$ values also indicate systematic loss of high-frequency detail
relative to sibling continuations of the same scenes. For
\texttt{noadaln}, this degradation should be considered together with its
weak action response and high relocation rate rather than as an isolated
sharpness failure.

\begin{table}[t]
\centering
\small
\caption{
    High-frequency degradation across the ablation grid. $\mathcal B$ is the
    excess sibling-relative loss of Laplacian sharpness between the base and
    end windows. The percentage column uses the fixed descriptive threshold
    $\mathcal B>150$. Dirty-source rollouts are excluded.
}
\label{tab:hf_losses}
\begin{tabular}{lcc}
\toprule
Variant & Median $\mathcal B$ & $\mathcal B>150$ \\
\midrule
16node  & $-191$ & $2\%$  \\
pca8    & $-184$ & $2\%$  \\
pca4    & $-161$ & $2\%$  \\
noatok  & $0$    & $17\%$ \\
pca2    & $+45$  & $20\%$ \\
4node   & $+144$ & $47\%$ \\
noadaln & $+203$ & $59\%$ \\
\bottomrule
\end{tabular}
\end{table}

Taken together, the ablations reveal several distinct quality failure modes.
Reduced PCA supervision primarily increases implausible content and geometric
corruption, reduced training scale increases relocation and progressive loss
of detail, and removing the AdaLN action pathway produces weak continuations
with severe relocation and high-frequency degradation. No single metric
therefore captures the full effect of the action-conditioning design.


\end{document}
